\documentclass[10pt]{article}
\usepackage[margin=0.9in]{geometry}
\usepackage{amsmath,amssymb}
\usepackage{booktabs,longtable}
\usepackage[hidelinks]{hyperref}
\title{Supplementary Material\\[2pt] \large Smaller 5-chromatic unit-distance graphs on two icosahedral spheres}
\author{Chandan Kumar Shukla\\ {\small Independent Researcher, India}}
\date{}
\begin{document}
\maketitle
\tableofcontents

\section{Overview}
\begin{center}\small
\begin{tabular}{lllrrl}
\toprule
graph & sphere & parent & vertices & edges & properties\\
\midrule
$H_{231}$ & $S^2(r_1)$, $r_1=\cos\frac{\pi}{10}$ & $G_{372}$ & 231 & 938 & $\chi=5$, vertex-critical, no Moser spindle\\
$H_{961}$ & $S^2(r_2)$, $r_2=\cos\frac{3\pi}{10}$ & $G_{972}$ & 961 & 4028 & $\chi=5$, vertex-critical, no Moser spindle\\
\bottomrule
\end{tabular}
\end{center}
$G_{372}$ and $G_{972}$ are the graphs of Voronov, Neopryatnaya and Dergachev (Discrete Math.\ 345 (2022)
113106). $H_{231}$ is the smallest graph found by our search; it is not claimed to be minimum. All files
named below are in the repository accompanying the paper (\url{https://github.com/chandanshukla1989/5chromatic-spheres}); \texttt{verify\_all.sh} runs every check.

\section{Exact coordinates and number fields}
For each sphere all coordinates lie in a number field $K$ of degree $8$ (Galois group of the Galois
closure $D_8$), a tower of quadratic extensions $\mathbb{Q}\subset\mathbb{Q}(\sqrt5)\subset
\mathbb{Q}(\sqrt5,b)\subset\mathbb{Q}(\sqrt5,b,g)=K$ with $b=\sqrt\beta$, $g=\sqrt\gamma$ (positive roots):
\begin{align*}
S^2(r_1):\quad & \beta=\tfrac72+\tfrac52\sqrt5\approx 9.0902,\quad
\gamma=\tfrac{15}{4}-\tfrac74\sqrt5+\bigl(\tfrac34-\tfrac14\sqrt5\bigr)b\approx 0.41269,\\
S^2(r_2):\quad & \beta=\tfrac{23}{2}+\tfrac{11}{2}\sqrt5\approx 23.798,\quad
\gamma=\tfrac54-\tfrac34\sqrt5+\bigl(-\tfrac14+\tfrac14\sqrt5\bigr)b\approx 1.0804.
\end{align*}
Every coordinate is written as $\sum q_j e_j$ with rational $q_j$ and
$e_j\in\{1,\sqrt5,b,\sqrt5 b,g,\sqrt5 g,bg,\sqrt5 bg\}$. The files
\texttt{H231\_exact\_coordinates.m} and \texttt{H961\_exact\_coordinates.m} give all vertices in this form
(Mathematica syntax, \texttt{Sqrt[...]}); the \texttt{.py} files give the same in Python syntax
(\texttt{(x)**(1/2)}; load with \texttt{sympy} in rational mode). All vertices of $H_{231}$ and the
eleven vertices deleted from $G_{972}$ are listed in Appendices~\ref{app:h231} and~\ref{app:removed}.

\section{Parent graphs and index mapping}
\texttt{construction/build\_G372.py} and \texttt{build\_G972.py} rebuild the parent graphs with $60$-digit
coordinates from the published minimal polynomials (and the closed form of the second coordinate of $w_8$).
The \emph{parent index} of a vertex is its position in the output: for $G_{372}$ the icosahedron (0--11),
then the orbits of $v_1$ (12--131), $v_2$ (132--251), $v_3$ (252--371); for $G_{972}$ the orbits of
$w_1,\dots,w_8$ (0--959, $120$ each), then the great icosahedron (960--971). Within an orbit the order is
that of a breadth-first enumeration of the group from its generators (as in the scripts). Column~2 of
\texttt{H231\_vertices.txt} and \texttt{H961\_vertices.txt} gives the parent index;
\texttt{verify\_graph.py} checks that each vertex coincides with the parent vertex of that index
(maximal difference $<10^{-30}$, the precision of the files). The group $\Gamma'$ used for $G_{972}$ is
generated, as in the authors' notebook, by $g_1,g_2,Mg_3M,g_4$; it equals $M\Gamma M$ ($M$: swap of $y$
and $z$), which we checked by enumerating both groups.

\section{Verification methodology}
\begin{center}\small
\begin{tabular}{p{0.32\textwidth}p{0.6\textwidth}}
\toprule
claim & method (program)\\
\midrule
vertices on the sphere, edges of length 1 & exact arithmetic in $K$ with PARI/GP (\texttt{verification/pari/exact372.gp}, \texttt{exact972.gp}); independently, exact arithmetic on the radical files with an own implementation of the tower (\texttt{check\_exact\_radicals.py}; perturbing a single coordinate was confirmed to be detected)\\
no missing unit edge & all non-adjacent pairs have $\bigl||p-q|-1\bigr|\ge %GAP231%$ ($H_{231}$), $\ge %GAP961%$ ($H_{961}$) with 30-digit coordinates (\texttt{verify\_graph.py})\\
no proper 4-colouring & DRAT proof checked by drat-trim; cross-checks with CaDiCaL, Glucose, MiniSat and a binary encoding (\texttt{crosscheck\_solvers.py}); SAT-free DSATUR search for $H_{231}$ (\texttt{dsatur.c})\\
proper 5-colouring & edge-by-edge check of \texttt{H*\_5colouring.txt} (\texttt{verify\_graph.py})\\
vertex-critical & edge-by-edge check of one 4-colouring of $G-v$ per vertex $v$ (\texttt{H*\_criticality\_certificates.txt}, \texttt{verify\_graph.py})\\
no Moser spindle & exhaustive search (\texttt{verify\_graph.py})\\
\bottomrule
\end{tabular}
\end{center}

\section{CNF encoding}
For a graph with vertices $0,\dots,n-1$ the variable $4v+c+1$ means ``vertex $v$ has colour $c$''
($c=0,\dots,3$). The files \texttt{H*\_4col.cnf} contain, for every vertex, the clause
$\bigvee_c x_{v,c}$ and the six clauses $\neg x_{v,c}\vee\neg x_{v,c'}$; for every edge $uv$ and colour
$c$ the clause $\neg x_{u,c}\vee\neg x_{v,c}$; and two unit clauses giving colours $0$ and $1$ to the
endpoints of the first edge (no loss of generality, by symmetry of colours). The formula is satisfiable if
and only if the graph is 4-colourable.

\section{DRAT verification}
Proofs were produced with \texttt{kissat --no-binary H231\_4col.cnf H231\_4col.drat} (Kissat 4.0.4) and
are stored compressed. To check:
\begin{verbatim}
xz -dk H231/H231_4col.drat.xz
drat-trim H231/H231_4col.cnf H231/H231_4col.drat -t 20000      # -> s VERIFIED
\end{verbatim}
and likewise for $H_{961}$ (about two minutes). Our drat-trim outputs are in
\texttt{H*\_drat-trim.log} and \texttt{results/drat-trim\_*.log}.

\section{5-colourings, criticality, Moser spindles}
\texttt{H*\_5colouring.txt} lists \emph{vertex colour} with colours $1$--$5$.
\texttt{H*\_criticality\_certificates.txt} has one line per vertex $v$: \emph{v string}, where character
$x$ of the string is the colour ($0$--$3$) of vertex $x$ in a proper 4-colouring of $G-v$ and
\texttt{-} for $x=v$. A Moser spindle is two rhombi (pairs of unit triangles sharing an edge) with a common
apex whose tips are adjacent; \texttt{verify\_graph.py} enumerates all of them (none found).

\section{Software versions}
Python 3.14.7; python-sat 1.9.dev15 (CaDiCaL 1.5.3, Glucose 4, MiniSat 2.2); mpmath 1.3.0; numpy 2.5.3;
sympy 1.14.0; PARI/GP 2.19.0 with galdata; Kissat 4.0.4; drat-trim (commit \texttt{2e3b2dc}, 2024-11-25);
Apple clang 21.0.0; XZ Utils 5.8.4; macOS (Darwin, arm64).

\section{Reproducing all checks}
\begin{verbatim}
pip install -r requirements.txt
DRAT_TRIM=/path/to/drat-trim ./verify_all.sh       # about 6 minutes on a laptop
\end{verbatim}
Individual steps and expected outputs are listed in \texttt{REPRODUCIBILITY.md}.

\section{File manifest (SHA-256)}
{\small
\begin{longtable}{p{0.45\textwidth}p{0.5\textwidth}}
\toprule
file & content\\
\midrule
%MANIFEST%
\bottomrule
\end{longtable}}

\appendix
\section{All vertices of $H_{231}$ in radicals}\label{app:h231}
Numbering as in Appendix A of the paper (vertex $i$ here is index $i-1$ in the edge list); parent index in
$G_{372}$, orbit, and colour in the proper 5-colouring. $b,g$ as in Section~2 for $S^2(r_1)$.

{\footnotesize\raggedright
%H231RADICALS%}

\section{The eleven vertices of $G_{972}$ deleted to obtain $H_{961}$}\label{app:removed}
Numbering as in Table 2 of the paper; $b,g$ as in Section~2 for $S^2(r_2)$. All other $961$ vertices of
$G_{972}$ form $H_{961}$; their exact coordinates are in \texttt{H961\_exact\_coordinates.m}.

{\footnotesize\raggedright
%REMOVED%}

\end{document}
